\subsection{The triangular two-object exclusion through fifteen arrows}

\begin{lemma}[The rational all-act obstruction at order ten]
\label{C:ten-rational-act-cones}
Let $|M|=10$ and let $S$ be any nonempty finite right $M$-set.
If the act cone has bounded rational homology, then all of its
rational degree-two products vanish.  In particular it cannot
have the rational cohomology ring of $\mathbb{CP}^2$.
\end{lemma}
\begin{proof}
A vertex with zero rational $H^2$ gives only relative degree-two
classes, whose products vanish.  Otherwise its minimal ideal has
at least two rows and columns.  At most two rows are covered by
Theorem~\ref{C:two-row-act-cups}, including nontrivial minimal
groups.  The order bound leaves only band shapes $3\times2$,
$4\times2$, and $3\times3$ when there are at least three rows.
For $3\times2$ there are four outside elements, and the rational
bounded-homology form of
Theorem~\ref{C:two-column-four-outside-act-cones} applies.
The other two shapes have at most two outside elements; by
Lemma~\ref{C:two-outside-act-obstruction}, bounded rational
homology forces their act cones to have zero fourth rational
homology.  This proves the assertion in every case.
\end{proof}

\begin{lemma}[Three states opposite a two-element right-zero band]
\label{C:three-state-right-zero-biset}
Let $N=\{1,p,q\}$, with $p^2=p$, $q^2=q$, $pq=q$, and $qp=p$.
Let $S$ be a three-element $(N,M)$-biset.  If its collage has
type $\mathbb{CP}^2$, then an ordinary right $M$-act cone on at
most three states has that type.  Consequently no such collage
with $|M|\leq17$ has type $\mathbb{CP}^2$.
\end{lemma}
\begin{proof}
Write $P,Q:S\to S$ for the two left actions.  The identities
$PQ=Q$ and $QP=P$ imply that both maps are retractions onto
the same nonempty subset $U$, on which both are the identity.
The right $M$-action commutes with both retractions.
If $|U|=1$, its point is fixed by both monoids.  The resulting
section $BM\times BN\to E_S$ of the edge double bar makes every
degree-two class on the collage relative to the vertex spaces.
Its square vanishes because the relative quotient is a suspension.
If $|U|=3$, the left action is trivial.  Since $BN$ is
contractible, collapsing this vertex gives an actual homotopy
equivalence with the ordinary three-state right-act cone over $M$.

Suppose $|U|=2$, and denote the remaining point by $z$.
If $Pz\ne Qz$, the left-action category is the union of the two
full subcategories on $\{z,Pz\}$ and $\{z,Qz\}$.  Each has
initial object $z$: its unique arrow from $z$ to its other object
is fixed under postcomposition by every $N$-endomorphism.  Their
intersection is the singleton $z$, and no simplex crosses between
the two arms.  Its nerve is therefore contractible.  Realizing
the right $M$-action bar now makes $E_S\to BM$ a homotopy
equivalence, by objectwise homotopy invariance of the bar construction.
The collage is contractible as well, since $BN$ is contractible.

If $Pz=Qz$, the common retraction $S\to U$ is an actual biset
retraction.  It gives a retraction from the collage $X_S$ to
$X_U$.  The latter is equivalent to the ordinary two-state
$M$-act cone and has zero positive-degree rational products by
Lemma~\ref{C:finite-line-act-cups}.  The inclusion of double bars
$E_U\subset E_S$ is a simplicial cofibration, and the quotient
of the two collage models is $\Sigma(E_S/E_U)$.  If
$\dim H^2(X_S;\mathbb Q)=1$, the retraction leaves two possibilities.
When $H^2(X_U;\mathbb Q)\ne0$, its pullback spans $H^2(X_S;\mathbb Q)$,
so the square vanishes.  When $H^2(X_U;\mathbb Q)=0$, every
degree-two class of $X_S$ is relative to $X_U$; the displayed
suspension quotient again makes its square zero.
These cases leave only the ordinary three-state act cone.
The last assertion follows from
Theorem~\ref{C:three-state-cone-minimum22}; the one- and two-state
cases have zero positive products directly.
\end{proof}

\begin{proposition}[The complete fifteen-arrow triangular calculation]
\label{C:fifteen-triangular}
No triangular finite category with two objects and at most fifteen
arrows has classifying space homotopy equivalent to $\mathbb{CP}^2$.
\end{proposition}
\begin{proof}
Proposition~\ref{C:fourteen-triangular} handles smaller totals.
Write $(m,n,s)$ for the two vertex orders and the number of cross
arrows.  A singleton biset gives a join, with zero positive
products.  For two cross arrows,
Lemma~\ref{C:two-state-biset-H2zero} permits restriction to
two vertices with nonzero rational $H^2$, leaving $(8,5,2)$
and $(7,6,2)$ up to opposition.
For $s\geq3$, some vertex has nonzero rational $H^2$; choose
it as the first vertex, taking the opposite category if necessary.

If the other vertex has order at most two, its rational
coinvariants are exact, and its actual orbit-act comparison
reduces the question to an ordinary act cone.  Vertices of order
at most nine are covered by the bounded-rational statement of
Lemma~\ref{C:nine-bounded-rank-one}.
The remaining profiles have a ten-element vertex and at most
four states, or an eleven-element vertex and at most three states.
The former are excluded by Lemma~\ref{C:ten-rational-act-cones},
which applies directly to the rational equivalence produced by
orbit coinvariants.  For the latter profile the opposite vertex
has order one, so this is an actual three-state act cone and
Theorem~\ref{C:three-state-cone-minimum22} applies.
One or two states were already excluded by the finite-line lemma.
Thus all residual profiles are precisely the seventeen rows below.
The larger order may occur in the second position when it is
the smaller vertex that supplies the nonzero degree-two class.

The first-vertex catalogs have $1,5,43,516$ tables at orders
five through eight, as in the preceding complete calculation.
At order nine, retain all $8166$ translation tables of the
$2\times2$, $2\times3$, and $3\times2$ minimal-band enumeration,
including its $124$ harmless rational-rank-zero entries.  Add
the identity-adjoined $2\times4$ and $4\times2$ bands and both
normalized $2\times2$ Rees ideals over $C_2$.  This gives a
complete $8170$-table cover of the possible positive-degree-two
vertices.  The last four tables are included literally: a
projective-dimension bound for an ordinary act cone alone would
not justify discarding them opposite a nonprojective corner.

For the second vertex take every ordinary class of orders
three through six.  At order seven, the independent complete
$17730$-class isomorphism-and-opposition census and its transposes
give $32823$ distinct literal tables.  This is an exhaustive
oriented cover; it is not asserted to be an ordinary isomorphism
class count.  In the two-state rows, use the positive-degree-two
catalogs for both vertices.  Any second vertex with projective
trivial right rational module is skipped only after supplying
and checking an actual vector $v$ satisfying
$vn=v$ for every $n$ and $\epsilon(v)=1$.
The orbit-act comparison and the bounded-rational statement of
Lemma~\ref{C:nine-bounded-rank-one} then apply.

The pointwise constraint solver assigns every value of every
right action and enforces all multiplication equations.  It
retains one representative under the full state-permutation
group.  For each such right action it assigns all left-action
values, imposes the opposite multiplication equations and every
commuting equation, and retains every solution.  No quotient by
a second, independently chosen state permutation is taken.
Thus every relative position of the two actions occurs.
All completed tables and action equations are checked exactly.
The disjoint normalized-bar outcomes are
\[
\begin{array}{c|r|r|r|r|r}
(m,n,s)&\text{bisets}&b_1\ne0&b_2\ne1&x^2=0&x^2\ne0\\\hline
(9,3,3)&817136&112228&479153&225333&422\\
(8,3,4)&480725&88987&222031&169460&247\\
(8,4,3)&366630&48500&207170&110750&210\\
(7,3,5)&358971&81252&182780&94735&204\\
(7,4,4)&333826&59897&180226&93477&226\\
(7,5,3)&271272&35609&154775&80384&504\\
(6,3,6)&368606&99538&149247&119786&35\\
(6,4,5)&433128&91248&240070&101810&0\\
(6,5,4)&355692&68244&181157&106003&288\\
(6,6,3)&324811&46065&174566&104030&150\\
(5,3,7)&921591&272192&482520&166879&0\\
(5,4,6)&1684292&318335&1071600&294357&0\\
(5,5,5)&1121618&233893&625012&262684&29\\
(5,6,4)&863223&174107&416858&271958&300\\
(5,7,3)&902825&131912&468763&301828&322\\
(8,5,2)&5604&516&1996&3084&8\\
(7,6,2)&2492&215&987&1199&91.
\end{array}
\]
Here every Betti number is over $\mathbb F_2$.
All $3036$ cases in the last column have nonzero third homology.
The other $2507757$ retained cases have explicit normalized
degree-three primitives for the square of a certified nonboundary
degree-two cocycle.  Consequently all $9612442$ bisets are excluded.
The first row also follows from
Lemma~\ref{C:three-state-right-zero-biset}: among the seven
three-element monoids, only the displayed right-zero extension
fails the checked rational right-trivial splitting test.

The catalogs and all projectivity witnesses are saved in
\path{code/outputs/b_fifteen_biset_input_manifest.json} and its
seventeen compressed input files.  The complete producer is
\path{code/b_fifteen_bisets.cpp}, with disjoint completion logs
and all retained certificates under
\path{code/outputs/b_fifteen_bisets_00.jsonl.gz} and
\path{code/outputs/b_fifteen_bisets_rest.jsonl.gz}.
The aggregate counts are in
\path{code/outputs/b_fifteen_bisets_summary.json}.
\end{proof}

\begin{remark}[Independent verification gate]
\label{C:fifteen-triangular-verification}
The independent pointwise action construction used here was
checked against every outcome of the fourteen-arrow enumeration.
For the new certificates, a separate normalized-bar implementation
rebuilds every literal category, checks all typed associativity
equations, verifies its second Betti number and the nonboundary
generator, and checks every claimed primitive on every degree-four
simplex.  The alternative certificates are checked by independent
third-homology ranks.  All seventeen profiles pass this independent
verification: $2510793$ retained certificates, $1859788417$ typed
associativity equations, $1167434995$ cocycle equations, and
$7023259440$ primitive equations.  Every one of the $3036$
alternative cases has independently recomputed positive third
homology.  The input preparation also checks all $31232$ saved
rational projectivity splittings with exact fractions.  The
separate final records are
\path{code/outputs/b_fifteen_biset_primitives_verified00.jsonl}
and
\path{code/outputs/b_fifteen_biset_primitives_verified_rest.jsonl}.
The verifier imports none of the producer's action enumeration,
cup solver, or cohomology routines.
\end{remark}

\begin{corollary}\label{C:fifteen-two-objects}
A finite category with at most two objects and classifying space
homotopy equivalent to $\mathbb{CP}^2$ has at least sixteen arrows.
\end{corollary}
\begin{proof}
For one object use Theorem~\ref{B:CP2-minimum23}.
A disconnected category cannot realize a connected space.
For two objects in one strongly connected component use
Theorem~\ref{C:twentytwo-strongly-connected}; the remaining
connected case is Proposition~\ref{C:fifteen-triangular}.
\end{proof}
